\documentclass[10pt,a4paper]{amsart}
\usepackage{amsmath,amssymb,amsthm,graphicx}
% Retained minimal preamble for the Stacks Project chapter more-algebra.tex (More on Algebra).
% The chapter file pinned-original-source.tex carries only \input{preamble} on its first line, so no
% preamble text lives inside the chapter itself.  The Stacks Project uses one single shared preamble for
% every chapter; the snapshot stages it identically as sources/stacks-chow/preamble.tex and
% sources/stacks-etale-cohomology/preamble.tex.  The block below is copied mechanically, line for line,
% from sources/stacks-chow/preamble.tex lines 14-18 (the Xy-pic setup) and 154-217 (the native theorem
% declarations together with the mathematical macros).  Nothing was added, removed or redefined; the
% wrapper class and the amsmath/amssymb/amsthm/graphicx base packages are supplied by the assembler.
% Chapter: more-algebra.tex -- More on Algebra.
% For commutative diagrams we use Xy-pic
\usepackage[all]{xy}

% We use 2cell for 2-commutative diagrams.
\xyoption{2cell}

% Theorem environments.
%
\theoremstyle{plain}
\newtheorem{theorem}[subsection]{Theorem}
\newtheorem{proposition}[subsection]{Proposition}
\newtheorem{lemma}[subsection]{Lemma}

\theoremstyle{definition}
\newtheorem{definition}[subsection]{Definition}
\newtheorem{example}[subsection]{Example}
\newtheorem{exercise}[subsection]{Exercise}
\newtheorem{situation}[subsection]{Situation}

\theoremstyle{remark}
\newtheorem{remark}[subsection]{Remark}
\newtheorem{remarks}[subsection]{Remarks}

\numberwithin{equation}{subsection}

% Macros
%
\def\lim{\mathop{\mathrm{lim}}\nolimits}
\def\colim{\mathop{\mathrm{colim}}\nolimits}
\def\Spec{\mathop{\mathrm{Spec}}}
\def\Hom{\mathop{\mathrm{Hom}}\nolimits}
\def\Ext{\mathop{\mathrm{Ext}}\nolimits}
\def\SheafHom{\mathop{\mathcal{H}\!\mathit{om}}\nolimits}
\def\SheafExt{\mathop{\mathcal{E}\!\mathit{xt}}\nolimits}
\def\Sch{\mathit{Sch}}
\def\Mor{\mathop{\mathrm{Mor}}\nolimits}
\def\Ob{\mathop{\mathrm{Ob}}\nolimits}
\def\Sh{\mathop{\mathit{Sh}}\nolimits}
\def\NL{\mathop{N\!L}\nolimits}
\def\CH{\mathop{\mathrm{CH}}\nolimits}
\def\proetale{{pro\text{-}\acute{e}tale}}
\def\etale{{\acute{e}tale}}
\def\QCoh{\mathit{QCoh}}
\def\Ker{\mathop{\mathrm{Ker}}}
\def\Im{\mathop{\mathrm{Im}}}
\def\Coker{\mathop{\mathrm{Coker}}}
\def\Coim{\mathop{\mathrm{Coim}}}

% Boxtimes
%
\DeclareMathSymbol{\boxtimes}{\mathbin}{AMSa}{"02}

%
% Macros for moduli stacks/spaces
%
\def\QCohstack{\mathcal{QC}\!\mathit{oh}}
\def\Cohstack{\mathcal{C}\!\mathit{oh}}
\def\Spacesstack{\mathcal{S}\!\mathit{paces}}
\def\Quotfunctor{\mathrm{Quot}}
\def\Hilbfunctor{\mathrm{Hilb}}
\def\Curvesstack{\mathcal{C}\!\mathit{urves}}
\def\Polarizedstack{\mathcal{P}\!\mathit{olarized}}
\def\Complexesstack{\mathcal{C}\!\mathit{omplexes}}
% \Pic is the operator that assigns to X its picard group, usage \Pic(X)
% \Picardstack_{X/B} denotes the Picard stack of X over B
% \Picardfunctor_{X/B} denotes the Picard functor of X over B
\def\Pic{\mathop{\mathrm{Pic}}\nolimits}
\def\Picardstack{\mathcal{P}\!\mathit{ic}}
\def\Picardfunctor{\mathrm{Pic}}
\def\Deformationcategory{\mathcal{D}\!\mathit{ef}}
\usepackage{cleveref}
\crefname{originallemma}{original lemma}{original lemmas}
\Crefname{originallemma}{Original lemma}{Original lemmas}
\makeatletter
\newlabel{more-algebra-lemma-finite-module-rank}{{0H9Z}{source}{}{sourcepointer-1}{}}
\newlabel{more-algebra-lemma-finite-module-rank@cref}{{[originallemma][0][]0H9Z}{[1][0][]source}{}{}{}}
\makeatother
\title{Four descriptions of a finite module over a domain relative to its rank}
\author{The Stacks Project Authors}
\date{stacks-project:more-algebra.tex:local-snapshot-20261012; GNU FDL 1.2 or later}
\begin{document}
\maketitle
\setcounter{subsection}{1}
\noindent\textit{Editorial note.} All delivered mathematical text is the author's own, copied from the pinned Stacks Project chapter file more-algebra.tex; no statement or proof step was rewritten, reordered or removed.\par\medskip
\noindent\textit{Editorial note.} The chapter file pinned-original-source.tex carries only \textbackslash{}input\{preamble\} on its first line, so no preamble text lives inside the chapter. The retained file sources/stacks-more-algebra/retained-preamble.tex is a mechanical, line-for-line copy of the project's single shared preamble as staged in the snapshot, namely sources/stacks-chow/preamble.tex lines 14-18 (Xy-pic setup) and 154-217 (theorem declarations and mathematical macros). Nothing was added, removed or redefined, no package was added, and the excerpt serves only as preamble, never as a span coordinate.\par\medskip
\noindent\textit{Editorial note.} Both counted spans start at the \textbackslash{}begin line of their environment, so the statement and the author's proof are delivered whole; surrounding prose of the section is deliberately outside the counted spans and is carried as standard prerequisites instead.\par\medskip
\noindent\textit{Editorial note.} The citation fields (source\_id, source\_version, authors, source\_url, license.name, source\_title) are identical to those of the cards already published from this chapter, so that the batch selector treats them as one source instead of deferring them as a same-source clash.\par\medskip
\noindent\textit{Allowed standard prerequisites (editorial).}\begin{itemize}
\item Rank over a domain, and the identification of M tensor K as a finite dimensional vector space for finite M. definition environment at source lines 5015-5019 of the same chapter. Supplies the r of the statement and the finite dimensional space in which the basis e\_1 to e\_r is chosen.
\item A finitely generated torsion module over a domain is annihilated by a single nonzero element. used twice at proof lines 5120-5121 and 5126-5128, for the cokernel of the comparison map and for the finite torsion module N. Produces g and g' whose product f makes the kernel and cokernel of the comparison map vanish after localisation.
\item A surjection onto a projective module splits. used at proof lines 5123-5124 to split M\_g as N plus R\_g\textasciicircum{}r. Isolates the finite torsion piece N whose annihilation gives the local freeness of M\_f.
\end{itemize}\medskip
\ifdefined\hypertarget\hypertarget{sourcepointer-1}{}\fi
\noindent\textit{Source reference.} Original lemma 0H9Z: focused-seeds/AGP/sources/stacks-snapshot-20261011/tags.txt, source lines 21544--21544. The lemma's own canonical tag, quoted so that readers can resolve the label inside the delivered text; it points at the counted statement itself and carries no proof content.\par\medskip
\begin{lemma}
\label{lemma-finite-module-rank}
Let $R$ be a domain. Let $M$ be a finite $R$-module. Then
\begin{enumerate}
\item $M$ has finite rank $r \geq 0$,
\item there is a map $M \to R^{\oplus r}$ whose kernel and cokernel
are torsion modules,
\item there is an $f \in R$, $f \not = 0$ such that $M_f$
is free of rank $r$, and
\item there is an injective map $R^{\oplus r} \to M$
whose cokernel is a torsion module.
\end{enumerate}
\end{lemma}

\begin{proof}
Since $M$ is finite, $M \otimes_R K$ is a finite dimensional
vector space over the fraction field $K$ of $R$. Thus the rank $r$ is finite.

\medskip\noindent
Choose generators $x_1, \ldots, x_n$ of $M$. Choose a basis
$e_1, \ldots, e_r$ for $M \otimes_R K$. We may write
$x_i \otimes 1 = \sum (a_{ij}/b_{ij}) e_j$ for  some
$a_{ij} , b_{ij} \in R$ with $b_{ij} \not = 0$. Setting
$b = \prod b_{ij}$ and replacing $a_{ij}$ by $a_{ij}b/b_{ij}$
we may assume that $x_i \otimes 1 = \sum (a_{ij}/b) e_j$.
Note that if we have a relation $\sum r_i x_i = 0$ in $M$
then $\sum r_ia_{ij} = 0$ in $R$ for all $j = 1, \ldots, r$.

\medskip\noindent
Proof of (2). The last comment above tells us the map
$M \to \bigoplus Re_j$ sending $x_i$ to $\sum a_{ij} e_j$.
is well defined. The fact that $e_1, \ldots, e_r$ forms a basis
of $M \otimes_R K$ tells us that this map becomes an
isomorphism upon tensoring by $K$. Hence this is a map as in (2).

\medskip\noindent
Proof of (3). The cokernel of the map $M \to R^{\oplus r}$
constructed in (2) is finite and torsion. Hence we can find
a $g \in R$, $g \not = 0$ which annihilates this cokernel.
In other words, such that the map $M_g \to R_g^{\oplus r}$
is surjective. Then we can choose a splitting
$M_g = N \oplus R_g^{\oplus r}$. Then $N$ is a finite
$R_g$-module with $N \otimes_{R_g} K = 0$. In other words,
$N$ is torsion and finite over $R_g$. Thus we can find
a $g' \in R$ which annihilates $N$. Setting $f= gg'$ we
see that (3) holds.

\medskip\noindent
To prove (4) choose a basis $y_1, \ldots, y_r$
of $M_f$ where $f$ is as in (3). Then write
$y_j = m_j/f^{t_j}$ for some $m_j \in M$ and
$t_j \geq 0$. The map $R^{\oplus r} \to M$ sending
the $j$th basis vector to $m_j$ will be as in (4).
Details omitted.
\end{proof}

\end{document}
